% Noether P08 Korean producer draft — UNCHECKED
% T05-U18; ED0002 lines6131--6138; source 620 LF B / 42546323B9EE39B3D2FD53955A6E33382E0305EDDD9D754F32A931062B2A7A6C
% Pointer v009 / ED0002: B06BE3530D9CF2E82B56FDBA7FE41D5D044DF2425DFA2A059D4939EAA2F7A6C2 / C9A125167ACB33D914EE4374B65AE7CDF0052F568371B8B77B720EA178ABF0E3
% Translation only; formulas are preserved for independent checking.

\emph{b) $\mathcal S$와 $\mathcal T$의 랭크를 비교하는 데에는 이에 대응하는 보조정리 b)가 쓰인다. $h(x,y,z)=0$이면 반드시 $g(\lambda;xy)=0$이다.}

이를 증명하기 위해, 3에 따라 $h(x,y,z)=0$이라는 것은 다음과 같은 각 미분연산자 거듭제곱의 곱이 0이 되는 것과 동치임을 생각하자.
\[
 \left(\frac{\partial}{\partial\lambda_1}\frac{\partial}{\partial x_1}
 +\frac{\partial}{\partial\lambda_2}\frac{\partial}{\partial y_1}\right)^{p_1}
 \left(\frac{\partial}{\partial\lambda_1}\frac{\partial}{\partial x_2}
 +\frac{\partial}{\partial\lambda_2}\frac{\partial}{\partial y_2}\right)^{p_2}g(\lambda;x,y),
 \qquad p_1+p_2=p.
\]
